\section{Boundary layers and increments}\label{sec:boundary}
A natural route to sums of squares at every order is induction on the
order. Define the increment in a common signed-symbol ring by
\begin{equation}\label{eq:increment}
                   R_n=(n-1)F_{n+1}-(n+1)F_n,\qquad n\ge2,
\end{equation}
so that
\[
 F_{n+1}=\frac{n+1}{n-1}F_n+\frac1{n-1}R_n .
\]
If every increment were SOS, every $F_n$ would be SOS. This section
shows that the increments eventually fail to be SOS. The reason is that
the stabilized corner has a simple discrete derivative, the
boundary-layer polynomial; the relation is an algebraic consequence of
depth weights and is independent of the proof of
Theorem~\ref{thm:finite-scale}.

In this section use one-based depths $1\le p,q\le \kappa$ and the same pure
and mixed wedge orientations as before. Put $d_{pq}=p+q-1$ and
\[
 \mathcal O_\delta=\sum_{p-q=\delta}z^{\rm mix}_{pq},\quad
 s_g=\sum_{p+q=g}z^{\rm mix}_{pq},\quad
 U_h=\sum_{p=1}^{\kappa-h}u_{p,p+h},\quad
 V_h=\sum_{p=1}^{\kappa-h}v_{p,p+h}.
\]
The boundary-layer polynomial is
\begin{equation}\label{eq:BL}
 \BL_\kappa=2\sum_{p,q=1}^\kappa d_{pq}(z^{\rm mix}_{pq})^2
       +2\sum_{p<q}d_{pq}(u_{pq}^2+v_{pq}^2)
       -2\sum_\delta \mathcal O_\delta^2.
\end{equation}
Its weights are linear in depth. The weights of $\calB_\kappa$ are products
$pq$; these are different polynomials.

Applying the Pl\"ucker relation to pairs of mixed wedges of the same
offset gives the exchange identity
\begin{equation}\label{eq:BL-exchange}
                 \sum_\delta \mathcal O_\delta^2=\sum_gs_g^2+2\sum_hU_hV_h.
\end{equation}
Each mixed antidiagonal pair occurs once, and each unordered pair of
pure wedges twice. This is the linear-weight counterpart of the
exchange that gives the matrices in Lemma~\ref{lem:corner}.

Set $\calB_0=0$. Let $\pi_j$ delete depths $1,\ldots,j-1$ and relabel the remaining ones
starting at $1$. Thus $\BL_{\kappa-j+1}\circ \pi_j$ uses only the tail of the
original variables. Direct summation of depth weights gives
\begin{equation}\label{eq:tail-decomposition}
 \calB_\kappa=\sum_{j=1}^\kappa\BL_{\kappa-j+1}\circ \pi_j,
 \qquad \calB_\kappa^{\rm tail}:=\sum_{j=2}^\kappa\BL_{\kappa-j+1}\circ \pi_j
             =\calB_{\kappa-1}\circ \pi_2.
\end{equation}
For example, the diagonal identity is
$\sum_{j=1}^{\min(p,q)}(p+q-2j+1)=pq$.
The offset kernel sums its common-depth shifts to
$\min(p,q,r,s)$. These prove every coefficient in
\eqref{eq:tail-decomposition}; also $\calB_\kappa=\calB_\kappa^{\rm tail}+\BL_\kappa$.

For the increment \eqref{eq:increment}, the variables at offsets $\pm n$ have zero weight in $F_n$.
For this step, retain the outer $\kappa$ depths
$\pm(n+1-p)$, $1\le p\le \kappa$, and call this restriction $\iota_{n,\kappa}$.
When $n\ge\max(2,2\kappa-1)$, direct shift multiplication gives
\begin{equation}\label{eq:increment-corner}
 F_{n+1}\circ\iota_{n,\kappa}=\calB_\kappa,\quad
 F_n\circ\iota_{n,\kappa}=\calB_\kappa^{\rm tail},\quad
 R_n\circ\iota_{n,\kappa}=(n-1)\BL_\kappa-2\calB_\kappa^{\rm tail}.
\end{equation}
The onset ensures that the two opposite boundary strings are disjoint.
The old weights in the restricted $F_n$ are $(p-1)(q-1)$ and the
minimum is $\min(p,q,r,s)-1$, exactly as the tail formula requires.

\begin{corollary}[Failure of the all-order increment condition]\label{cor:boundary}
There is a least integer $\kappa_*$ for which $\BL_{\kappa_*}$ is not SOS.
For every $\kappa\ge \kappa_*$, $\BL_\kappa$ is not SOS, and for every
$n\ge\max(2,2\kappa_*-1)$, $R_n$ is not SOS.
\end{corollary}
\begin{proof}
If every $\BL_\kappa$ were SOS, \eqref{eq:tail-decomposition} would make
every $\calB_\kappa$ SOS, contradicting Corollary~\ref{cor:threshold}.
Restriction to the first $\kappa_*$ depths sends $\BL_\kappa$ to $\BL_{\kappa_*}$,
which proves the all-larger-width assertion.
By minimality, every summand in $\calB_{\kappa_*}^{\rm tail}$ is SOS. If $R_n$ were SOS at
one of the stated orders, \eqref{eq:increment-corner} would give
\[
 (n-1)\BL_{\kappa_*}=R_n\circ\iota_{n,\kappa_*}+2\calB_{\kappa_*}^{\rm tail}
\]
as a sum of squares, a contradiction.
\end{proof}
